\subsection{Objective 5: extended runtime benchmark material}
\label{sec:supp-runtime-benchmark}

\subsubsection{Print scope, export format, and HRU scaling ladder}
\label{sec:supp-runtime-print-scope}

Print-scope savings dominate calibration-oriented runs (Table~\ref{tab:runtime-benchmark-print-scope}; Figure~\ref{fig:runtime-benchmark-print-scope}).
NetCDF is faster than formatted text under full export (Table~\ref{tab:runtime-benchmark-nc-vs-txt}).
Production compiler rows are in Table~\ref{tab:runtime-benchmark-compiler}; the six-point HRU scaling ladder is in Table~\ref{tab:runtime-benchmark-hru-scaling}.

\begin{table}[htbp]
  \centering
  \caption{Print-scope comparison (full vs calibration filtered NetCDF).}
  \label{tab:runtime-benchmark-print-scope}
  \small
  \begin{tabular*}{\textwidth}{@{\extracolsep{\fill}}lllrrr@{}}
\toprule
Size class & Model & Print scope & Wall (s) & s/day & channel\_sd (MB) \\
\midrule
S & \texttt{03080102} & Full export & 21.13 & 0.05 & 5.46 \\
S & \texttt{03080102} & Calibration filtered & 5.42 & 0.01 & 0.45 \\
M & \texttt{09471300} & Full export & 269.37 & 0.68 & 175.26 \\
M & \texttt{09471300} & Calibration filtered & 99.79 & 0.26 & 0.56 \\
L & \texttt{03100101} & Calibration filtered & 1378.13 & 3.37 & 4.89 \\
\bottomrule
\end{tabular*}
%
\end{table}

\begin{figure}[htbp]
  \centering
  \includegraphics[width=0.72\textwidth]{final/fig-runtime-benchmark-print-scope.png}
  \caption{Print-scope effect for the small and medium benchmark domains.}
  \label{fig:runtime-benchmark-print-scope}
\end{figure}

\begin{table}[htbp]
  \centering
  \caption{Full-export NetCDF vs formatted text (one-year run).}
  \label{tab:runtime-benchmark-nc-vs-txt}
  \small
  \begin{tabular*}{\textwidth}{@{\extracolsep{\fill}}lllrrr@{}}
\toprule
Size class & Model & Format & Wall (s) & s/day & Total output (GB) \\
\midrule
S & \texttt{03080102} & NC & 21.13 & 0.05 & 0.13 \\
S & \texttt{03080102} & TXT & 26.15 & 0.07 & 0.61 \\
M & \texttt{09471300} & NC & 269.37 & 0.68 & 3.07 \\
M & \texttt{09471300} & TXT & 596.98 & 1.6 & 14.57 \\
\bottomrule
\end{tabular*}
%
\end{table}

\begin{table}[htbp]
  \centering
  \caption{Production ifx \texttt{-O3 -ipo} on filtered profile.}
  \label{tab:runtime-benchmark-compiler}
  \small
  \begin{tabular*}{\textwidth}{@{\extracolsep{\fill}}lllrr@{}}
\toprule
Size class & Model & Build & Wall (s) & s/day \\
\midrule
S & \texttt{03080102} & ifx -O3 -ipo & 5.05 & 0.01 \\
M & \texttt{09471300} & ifx -O3 -ipo & 89.55 & 0.23 \\
L & \texttt{03100101} & ifx -O3 -ipo & 1352.57 & 3.29 \\
\bottomrule
\end{tabular*}
%
\end{table}

\begin{table}[htbp]
  \centering
  \caption{Six-point simulation benchmark ladder (S/M/L + X20/X40/X60 runtime scaling cases; filtered profile; production ifx).}
  \label{tab:runtime-benchmark-hru-scaling}
  \small
  \begin{tabular*}{\textwidth}{@{\extracolsep{\fill}}lllrrrr@{}}
\toprule
Size class & Model & HRUs & Channels & Wall (s) & s/day & Init (s) \\
\midrule
S & \texttt{03080102} & 473 & 45 & 5.05 & 0.01 & 1.0 \\
M & \texttt{09471300} & 11284 & 350 & 89.55 & 0.23 & 5.0 \\
X20 & \texttt{03152000} & 19530 & 1615 & 156.68 & 0.39 & 13.0 \\
X40 & \texttt{07174000} & 36855 & 2370 & 297.37 & 0.75 & 22.0 \\
X60 & \texttt{15060105} & 51685 & 17296 & 1436.92 & 3.77 & 56.0 \\
L & \texttt{03100101} & 94303 & 8181 & 1352.57 & 3.29 & 149.0 \\
\bottomrule
\end{tabular*}
%
\end{table}

\subsubsection{Full compiler matrix (Scenario 3)}
\label{sec:supp-runtime-compiler}

\begin{table}[htbp]
  \centering
  \caption{Supplementary Table~S2. Compiler optimization matrix on the filtered calibration print profile (365-day runs). ``Prod.'' marks the production ifx \texttt{-O3 -ipo} build used elsewhere.}
  \label{tab:supp-runtime-compiler}
  \scriptsize
  \setlength{\tabcolsep}{3pt}
  \renewcommand{\arraystretch}{1.05}
  \begin{tabular*}{\textwidth}{@{\extracolsep{\fill}}lllrrrl@{}}
\toprule
Size class & Model & Build & Wall (s) & s/day & vs ref (\%) & Prod. \\
\midrule
S & \texttt{03080102} & gfortran -O2 & 5.54 & 0.01 & -1.2 &  \\
S & \texttt{03080102} & gfortran -O3 & 6.5 & 0.01 & 15.9 &  \\
S & \texttt{03080102} & ifx -O2 & 5.61 & 0.01 & 0.0 &  \\
S & \texttt{03080102} & ifx -O3 & 5.49 & 0.01 & -2.1 &  \\
S & \texttt{03080102} & ifx -O3 -xHost & 6.04 & 0.01 & 7.7 &  \\
S & \texttt{03080102} & ifx -O3 -ipo & 5.05 & 0.01 & -10.0 & yes \\
M & \texttt{09471300} & gfortran -O2 & 111.56 & 0.28 & 15.3 &  \\
M & \texttt{09471300} & gfortran -O3 & 156.99 & 0.41 & 62.2 &  \\
M & \texttt{09471300} & ifx -O2 & 96.79 & 0.25 & 0.0 &  \\
M & \texttt{09471300} & ifx -O3 & 92.2 & 0.24 & -4.7 &  \\
M & \texttt{09471300} & ifx -O3 -xHost & 104.72 & 0.27 & 8.2 &  \\
M & \texttt{09471300} & ifx -O3 -ipo & 89.55 & 0.23 & -7.5 & yes \\
L & \texttt{03100101} & gfortran -O2 & 1269.09 & 2.93 & -8.9 &  \\
L & \texttt{03100101} & ifx -O2 & 1393.06 & 3.39 & 0.0 &  \\
L & \texttt{03100101} & ifx -O3 & 1391.73 & 3.38 & -0.1 &  \\
L & \texttt{03100101} & ifx -O3 -xHost & 1363.69 & 3.33 & -2.1 &  \\
L & \texttt{03100101} & ifx -O3 -ipo & 1352.57 & 3.29 & -2.9 & yes \\
\bottomrule
\end{tabular*}
%
\end{table}

\subsubsection{VTune CPU splits (Scenario 4, large basins)}
\label{sec:supp-runtime-vtune}

\begin{table}[htbp]
  \centering
  \caption{Supplementary Table~S3. Intel VTune hotspot shares on filtered calibration runs (90 simulated days). Daily-loop percentages are shown for the repeating simulation phase; Init is reported separately.}
  \label{tab:supp-runtime-vtune}
  \scriptsize
  \setlength{\tabcolsep}{4pt}
  \renewcommand{\arraystretch}{1.08}
  \begin{tabular*}{\textwidth}{@{\extracolsep{\fill}}llrrrrrr@{}}
\toprule
Tier & Model & Days & HRU (\%) & Channel (\%) & strcmp (\%) & memset (\%) & Init (\%) \\
\midrule
L & \texttt{03100101} & 90 & 16.6 & 11.8 & 20.7 & 18.7 & 5.1 \\
X60 & \texttt{15060105} & 90 & 8.0 & 47.6 & 5.6 & 6.3 & 1.8 \\
\bottomrule
\end{tabular*}
%
\end{table}

Peace River allocates a larger share to HRU processing; Upper Gila allocates nearly half of daily-loop CPU to channel routing despite fewer HRUs, consistent with the six-point ladder inversion discussed in Section~\ref{sec:results-runtime-benchmark}.

\subsubsection{Calibration wall-time extrapolation}
\label{sec:supp-runtime-calibration-cost}

Scenario~4 measures one filtered forward run per model (365 simulated days).
A full particle-swarm calibration repeats comparable runs while searching parameter space.
Typical watershed calibrations adjust on the order of 30--40 parameters; our production setup often uses 24--48 particles and a maximum of about 75 iterations, with early stop when the objective stagnates (often between 30 and 75 iterations).

We extrapolate each model's measured initialization and daily-loop split to a seven-year simulation window and estimate calibration wall time as
\[
\text{wall} \approx N_{\mathrm{iter}} \times \left\lceil N_{\mathrm{particles}} / N_{\mathrm{cores}} \right\rceil \times \bigl(t_{\mathrm{init}} + 7 \times t_{\mathrm{daily}}\bigr),
\]
using Scenario~4 timing and the filtered calibration NetCDF profile.
Reported ranges span 24--48 particles, 30--75 iterations, and example servers with 16, 32, or 64 cores.
RAM limits are ignored in the formula; in practice memory can cap concurrent particles even when core count is higher.
Doubling cores helps only until $N_{\mathrm{cores}} \geq N_{\mathrm{particles}}$ each iteration; further speedup requires fewer particles, fewer iterations, a shorter scored window, or the print-scope reductions in Scenarios~1--2.

\subsubsection{Production calibration optimizations beyond forward-run timing}
\label{sec:supp-runtime-optimization-stack}

The measured scenarios isolate single forward runs.
SWATGenX production calibration additionally applies:
\begin{itemize}
  \item \textbf{Gauge-limited printing} --- \texttt{print\_filter.prt} restricts daily \texttt{channel\_sd} to calibration gauges (\texttt{swift\_out=0}, \texttt{wq\_cha=0}).
  \item \textbf{NetCDF export} --- \texttt{cdfout=y} when full daily/monthly exports are required (fork extension; Table~\ref{tab:runtime-benchmark-nc-vs-txt}).
  \item \textbf{Compiler selection} --- Intel ifx \texttt{-O3 -ipo} on the reference host (Table~\ref{tab:runtime-benchmark-compiler}; Table~S2).
  \item \textbf{Shared climate I/O} --- common weather paths via \texttt{file.cio} so parallel particles reuse the same forcing files where possible.
  \item \textbf{Optional coarser landuse/soil grids} at model build to reduce HRU cardinality before calibration.
  \item \textbf{PSO pool sizing and stagnation stop} --- limit particles and terminate when the objective plateaus.
\end{itemize}
These workflow choices reduce batch overhead beyond what a one-year forward-run ladder alone captures.
